\frfilename{mt223.tex}
\versiondate{9.9.04}
\copyrightdate{1995}

\def\chaptername{Teorema Dasar Kalkulus}
\def\sectionname{Teorema-teorema kerapatan Lebesgue}
\def\undpsi{\underline{\psi}}

\newsection{223}

Sekarang saya beralih ke sekelompok hasil yang dapat dipandang sebagai
korolari dari Teorema 222E, tetapi juga berkembang pesat sebagai hasil-hasil tersendiri,
termasuk kemungkinan generalisasi penting yang akan dibahas dalam Bab 26.
Gagasannya ialah bahwa setiap fungsi terukur $f$ pada $\Bbb R$ bersifat
`kontinu' hampir di mana-mana dalam berbagai pengertian yang sangat lemah;
untuk hampir setiap $x$, nilai $f(x)$ ditentukan oleh perilaku $f$ di dekat
$x$, dalam arti bahwa $f(y)\bumpeq f(x)$ untuk `sebagian besar' $y$ di dekat
$x$. Mungkin perlu saya katakan bahwa, walaupun saya menyarankan bagian ini
sebagai persiapan untuk Bab 26 dan juga mengandalkannya dalam Bab 28, saya
tidak akan merujuknya lagi dalam bab ini. Karena itu, pembaca yang ingin
segera mencirikan integral tak tentu dapat langsung menuju \S224.

\leader{223A}{Teorema Kerapatan Lebesgue: bentuk integral} Misalkan $I$
suatu interval dalam $\Bbb R$, dan $f$ suatu fungsi bernilai real yang
terintegralkan pada $I$. Maka

$$\eqalign{f(x)
&=\lim_{h\downarrow 0}\Bover1{h}\int_{x}^{x+h}f
=\lim_{h\downarrow 0}\Bover1h\int_{x-h}^xf
=\lim_{h\downarrow 0}\Bover1{2h}\int_{x-h}^{x+h}f\cr}$$

\noindent untuk hampir setiap $x\in I$.


\proof{ Tetapkan $F(x)=\int_{I\cap\ooint{-\infty,x}}f$. Dari 222G kita
mengetahui bahwa

$$\eqalign{f(x)=F'(x)
&=\lim_{h\downarrow 0}\Bover1{h}(F(x+h)-F(x))
   =\lim_{h\downarrow 0}\Bover1{h}\int_{x}^{x+h}f\cr
&=\lim_{h\downarrow 0}\Bover1h(F(x)-F(x-h))
   =\lim_{h\downarrow 0}\Bover1h\int_{x-h}^xf\cr
&=\lim_{h\downarrow 0}\Bover1{2h}(F(x+h)-F(x-h))
   =\lim_{h\downarrow 0}\Bover1{2h}\int_{x-h}^{x+h}f\cr}$$

\noindent untuk hampir setiap $x\in I$.
}%end of proof of 223A

\leader{223B}{Korolari} Misalkan $E\subseteq\Bbb R$ suatu himpunan terukur.
Maka

\Centerline{$\lim_{h\downarrow 0}\Bover1{2h}\mu(E\cap[x-h,x+h])=1$
untuk hampir setiap $x\in E$,}

\Centerline{$\lim_{h\downarrow 0}\Bover1{2h}\mu(E\cap[x-h,x+h])=0$
untuk hampir setiap $x\in \Bbb R\setminus E$.}


\proof{ Ambil $n\in \Bbb N$. Dengan menerapkan 223A pada
$f=\chi(E\cap[-n,n])$, kita melihat bahwa

\Centerline{$\lim_{h\downarrow 0}\Bover1{2h}\int_{x-h}^{x+h}f
=\lim_{h\downarrow 0}\Bover1{2h}\mu(E\cap[x-h,x+h])$}

\noindent kapan pun $x\in\ooint{-n,n}$ dan salah satu limit tersebut ada,
sehingga

\Centerline{$\lim_{h\downarrow 0}\Bover1{2h}\mu(E\cap[x-h,x+h])=1$
untuk hampir setiap $x\in E\cap[-n,n]$,}

\Centerline{$\lim_{h\downarrow 0}\Bover1{2h}\mu(E\cap[x-h,x+h])=0$
untuk hampir setiap $x\in [-n,n]\setminus E$.}

\noindent Karena $n$ sebarang, hasilnya terbukti.
}%end of proof of 223B

\cmmnt{\medskip

\noindent{\bf Catatan} Untuk suatu himpunan terukur $E\subseteq\Bbb R$,
suatu titik $x$ sedemikian sehingga
$\lim_{h\downarrow 0}\Bover1{2h}\mu(E\cap[x-h,x+h])=1$ kadang-kadang
disebut suatu {\bf titik kerapatan} dari $E$.
}%end of comment

\leader{223C}{Korolari} Misalkan $f$ suatu fungsi bernilai real terukur yang
didefinisikan hampir di mana-mana dalam $\Bbb R$. Maka, untuk hampir setiap
$x\in\Bbb R$,

\Centerline{$\lim_{h\downarrow 0}\Bover1{2h}\mu\{y:y\in\dom f,\,|y-x|\le
h,\,|f(y)-f(x)|\le\epsilon\}=1$,}

\Centerline{$\lim_{h\downarrow 0}\Bover1{2h}\mu\{y:y\in\dom f,\,|y-x|\le
h,\,|f(y)-f(x)|\ge\epsilon\}=0$}

\noindent untuk setiap $\epsilon>0$.

\proof{ Untuk $q$, $q'\in\Bbb Q$, tetapkan

\Centerline{$D_{qq'}=\{x:x\in\dom f,\,q\le f(x)<q'\}$,}

\noindent sehingga $D_{qq'}$ terukur, dan tetapkan

\Centerline{$C_{qq'}=\{x:x\in D_{qq'},\,\lim_{h\downarrow 0}
\Bover1{2h}\mu(D_{qq'}\cap [x-h,x+h])=1\}$,}


\noindent sehingga $D_{qq'}\setminus C_{qq'}$ terabaikan, menurut 223B;
sekarang tetapkan

\Centerline{$C=\dom f\setminus\bigcup_{q,q'\in\Bbb Q}(D_{qq'}\setminus
C_{qq'})$,}

\noindent sehingga $C$ koterabaikan. Jika $x\in C$ dan $\epsilon>0$, maka
terdapat $q$, $q'\in \Bbb Q$ sedemikian sehingga
$f(x)-\epsilon\le q\le f(x)<q'\le f(x)+\epsilon$. Jadi $x$ termasuk dalam
$D_{qq'}$ dan karena itu dalam $C_{qq'}$, dan sekarang

$$\eqalign{\liminf_{h\downarrow 0}\Bover1{2h}
&\mu\{y:y\in \dom f\cap [x-h,x+h],\,|f(y)-f(x)|\le\epsilon\}\cr
&\ge\liminf_{h\downarrow 0}
\Bover1{2h}\mu(D_{qq'}\cap [x-h,x+h])
=1,\cr}$$

\noindent sehingga

\Centerline{$\lim_{h\downarrow 0}
\Bover1{2h}\mu\{y:y\in \dom f\cap [x-h,x+h],\,|f(y)-f(x)|\le\epsilon\}
=1$.}

Dengan segera diperoleh bahwa

\Centerline{$\lim_{h\downarrow 0}
\Bover1{2h}\mu\{y:y\in \dom f\cap [x-h,x+h],\,|f(y)-f(x)|>\epsilon\}
=0$}

\noindent untuk hampir setiap $x$; tetapi karena $\epsilon$ sebarang, hal ini
juga benar untuk $\bover12\epsilon$, sehingga sebenarnya

\Centerline{$\lim_{h\downarrow 0}
\Bover1{2h}\mu\{y:y\in \dom f\cap [x-h,x+h],\,|f(y)-f(x)|\ge\epsilon\}
=0$}

\noindent untuk hampir setiap $x$.
}%end of proof of 223C

\leader{223D}{Teorema} Misalkan $I$ suatu interval dalam $\Bbb R$, dan $f$
suatu fungsi bernilai real yang terintegralkan pada $I$. Maka

\Centerline{$\lim_{h\downarrow 0}\Bover1{2h}
\int_{x-h}^{x+h}|f(y)-f(x)|dy=0$}

\noindent untuk hampir setiap $x\in I$.

\proof{{\bf (a)} Mula-mula andaikan bahwa $I$ suatu interval terbuka
terbatas $\ooint{a,b}$. Untuk setiap $q\in\Bbb Q$, tetapkan
$g_q(x)=|f(x)-q|$ untuk $x\in I\cap\dom f$; maka $g_q$ terintegralkan pada
$I$, dan

\Centerline{$\lim_{h\downarrow 0}\Bover1{2h}
\int_{x-h}^{x+h}g_q=g_q(x)$}

\noindent untuk hampir setiap $x\in I$, menurut 223A. Dengan menetapkan

\Centerline{$E_q=\{x:x\in I\cap\dom f,\,\lim_{h\downarrow 0}
\Bover1{2h}\int_{x-h}^{x+h}g_q=g_q(x)\}$,}

\noindent kita memperoleh bahwa $I\setminus E_q$ terabaikan, sehingga
$I\setminus E$ terabaikan, dengan $E=\bigcap_{q\in\Bbb Q}E_q$. Sekarang

\Centerline{$\lim_{h\downarrow 0}\Bover1{2h}
\int_{x-h}^{x+h}|f(y)-f(x)|dy=0$}

\noindent untuk setiap $x\in E$. \Prf\ Ambil $x\in E$ dan $\epsilon>0$.
Maka terdapat $q\in\Bbb Q$ sedemikian sehingga
$|f(x)-q|\le\epsilon$, sehingga

\Centerline{$|f(y)-f(x)|\le|f(y)-q|+\epsilon=g_q(y)+\epsilon$}

\noindent untuk setiap $y\in I\cap\dom f$, dan

$$\eqalign{\limsup_{h\downarrow 0}\Bover1{2h}
\int_{x-h}^{x+h}|f(y)-f(x)|dy
&\le\limsup_{h\downarrow 0}\Bover1{2h}
\int_{x-h}^{x+h}g_q(y)+\epsilon\,dy\cr
&=\epsilon+g_q(x)
\le 2\epsilon.\cr}$$

\noindent Karena $\epsilon$ sebarang,

\Centerline{$\lim_{h\downarrow 0}\Bover1{2h}
\int_{x-h}^{x+h}|f(y)-f(x)|dy=0$,}

\noindent sebagaimana diperlukan.\ \Qed

\medskip

{\bf (b)} Jika $I$ suatu interval terbuka tak terbatas, terapkan (a) pada
interval-interval $I_n=I\cap\ooint{-n,n}$ untuk melihat bahwa limit tersebut
nol hampir di mana-mana dalam setiap $I_n$, dan karena itu pada $I$. Jika $I$
suatu interval sebarang, perhatikan bahwa interval itu berbeda dari suatu
interval terbuka pada paling banyak dua titik. Karena kita hanya mencari sifat
yang berlaku hampir di mana-mana, titik-titik ini dapat diabaikan.
}%end of proof of 223D

\medskip

\noindent{\bf Catatan} Himpunan

\Centerline{$\{x:x\in\dom f,\,\lim_{h\downarrow 0}\Bover1{2h}
\int_{x-h}^{x+h}|f(y)-f(x)|dy=0\}$}

\noindent kadang-kadang disebut {\bf himpunan Lebesgue} dari $f$.

\leader{223E}{Fungsi bernilai kompleks}\cmmnt{ Hasil-hasil di atas telah
saya nyatakan dalam bahasa fungsi bernilai real, karena itulah wahana yang
paling alami bagi gagasan-gagasannya. Namun terdapat penerapan yang sangat
penting dengan fungsi-fungsi bernilai kompleks, jadi saya tuliskan secara
eksplisit pernyataan-pernyataan yang bersangkutan di sini. Dalam setiap kasus,
buktinya elementer: cukup menerapkan hasil yang bersesuaian (223A, 223C, atau
223D) pada bagian real dan bagian imajiner fungsi $f$.}

\header{223Ea}{\bf (a)} Misalkan $I$ suatu interval dalam $\Bbb R$, dan $f$
suatu fungsi bernilai kompleks yang terintegralkan pada $I$. Maka

$$\eqalign{f(x)
=\lim_{h\downarrow 0}\Bover1{h}\int_{x}^{x+h}f
=\lim_{h\downarrow 0}\Bover1h\int_{x-h}^xf
=\lim_{h\downarrow 0}\Bover1{2h}\int_{x-h}^{x+h}f\cr}$$

\noindent untuk hampir setiap $x\in I$.

\header{223Eb}{\bf (b)} Misalkan $f$ suatu fungsi bernilai kompleks terukur
yang didefinisikan hampir di mana-mana dalam $\Bbb R$. Maka, untuk hampir
setiap $x\in\Bbb R$,

\Centerline{$\lim_{h\downarrow 0}\Bover1{2h}\mu\{y:y\in\dom f,\,|y-x|\le
h,\,|f(y)-f(x)|\ge\epsilon\}=0$}

\noindent untuk setiap $\epsilon>0$.

\header{223Ec}{\bf (c)} Misalkan $I$ suatu interval dalam $\Bbb R$, dan $f$
suatu fungsi bernilai kompleks yang terintegralkan pada $I$. Maka

\Centerline{$\lim_{h\downarrow 0}\Bover1{2h}
\int_{x-h}^{x+h}|f(y)-f(x)|dy=0$}

\noindent untuk hampir setiap $x\in I$.

\exercises{
\leader{223X}{Latihan dasar $\pmb{>}$(a)} Misalkan $E\subseteq[0,1]$ suatu
himpunan terukur sedemikian sehingga terdapat $\alpha>0$ dengan
$\mu(E\cap[a,b])\ge\alpha(b-a)$ kapan pun $0\le a\le b\le 1$.
Tunjukkan bahwa $\mu E=1$.
%223B

\spheader 223Xb Misalkan $A\subseteq\Bbb R$ sembarang himpunan. Tunjukkan
bahwa $\lim_{h\downarrow 0}\Bover1{2h}\mu^*(A\cap[x-h,x+h])=1$ untuk
hampir setiap $x\in A$. \Hint{terapkan 223B pada suatu selubung terukur
$E$ dari $A$.}
%223B

\spheader 223Xc Misalkan $E$, $F\subseteq\Bbb R$ himpunan-himpunan terukur,
dan misalkan $x\in\Bbb R$ merupakan titik kerapatan dari keduanya.
Tunjukkan bahwa $x$ merupakan titik kerapatan dari $E\cap F$.
%223B

\spheader 223Xd Misalkan $E\subseteq\Bbb R$ suatu himpunan terukur tak
terabaikan. Tunjukkan bahwa untuk setiap $n\in\Bbb N$ terdapat
$\delta>0$ sedemikian sehingga $\bigcap_{i\le n}E+x_i$ tak kosong kapan pun
$x_0,\ldots,x_n\in\Bbb R$ memenuhi $|x_i-x_j|\le\delta$ untuk semua $i$,
$j\le n$. \Hint{carilah suatu interval tak-sepele $I$ sedemikian sehingga
$\mu(E\cap I)>\bover{n}{n+1}\mu I$.}
%223B

\spheader 223Xe Misalkan $f$ sembarang fungsi bernilai real yang didefinisikan
hampir di mana-mana dalam $\Bbb R$. Tunjukkan bahwa $\lim_{h\downarrow
0}\Bover1{2h}\mu^*\{y:y\in\dom f,\,|y-x|\le
h,\,|f(y)-f(x)|\le\epsilon\}=1$ untuk hampir setiap $x\in\Bbb R$.
\Hint{gunakan argumen 223C, tetapi dengan 223Xb sebagai pengganti 223B.}
%223C

\sqheader 223Xf Misalkan $I$ suatu interval dalam $\Bbb R$, dan $f$
suatu fungsi bernilai real yang terintegralkan pada $I$. Tunjukkan bahwa
$\lim_{h\downarrow 0}\bover1{h}
\int_{x}^{x+h}|f(y)-f(x)|dy=0$ untuk hampir setiap $x\in I$.
%223D

\spheader 223Xg Misalkan $E$, $F\subseteq\Bbb R$ himpunan-himpunan terukur,
dan andaikan bahwa $F$ terbatas serta berukuran tak nol. Misalkan
$x\in\Bbb R$ sedemikian sehingga
$\lim_{h\downarrow 0}\Bover1{2h}\mu(E\cap[x-h,x+h])=1$. Tunjukkan bahwa
$\lim_{h\downarrow 0}\Bover{\mu(E\cap(x+hF))}{h\,\mu F}=1$.
({\it Petunjuk:\/} berguna untuk mengetahui bahwa $\mu(hF)=h\mu F$
(134Ya, 263A). Tunjukkan bahwa jika $F\subseteq[-M,M]$, maka

\Centerline{$\Bover1{2hM}\mu(E\cap[x-hM,x+hM])
\le 1-\Bover{\mu F}{2M}
  \bigl(1-\Bover{\mu(E\cap(x+hF))}{h\,\mu F}\bigr)$.)}

\noindent (Bandingkan 223Ya.)
%223B

\spheader 223Xh Misalkan $f$ suatu fungsi bernilai real yang terintegralkan
pada $\Bbb R$, dan $E$ himpunan Lebesgue dari $f$. Tunjukkan bahwa
$\lim_{h\downarrow 0}\bover1{2h}\int_{x-h}^{x+h}|f(t)-c|dt=|f(x)-c|$
untuk setiap $x\in E$ dan $c\in\Bbb R$.
%223D

\spheader 223Xi Misalkan $f$ suatu fungsi bernilai real terintegralkan yang
didefinisikan hampir di mana-mana dalam $\Bbb R$. Misalkan $x\in\dom f$
sedemikian sehingga
$\lim_{n\to\infty}\Bover{n}2\int^{x+1/n}_{x-1/n}|f(y)-f(x)|=0$.
Tunjukkan bahwa $x$ termasuk dalam himpunan Lebesgue dari $f$.
%223D

\spheader 223Xj Misalkan $f$ suatu fungsi bernilai real terintegralkan yang
didefinisikan hampir di mana-mana dalam $\Bbb R$, dan $x$ sembarang titik
dari himpunan Lebesgue dari $f$. Tunjukkan bahwa untuk setiap $\epsilon>0$
terdapat $\delta>0$ sedemikian sehingga, kapan pun $I$ suatu interval
tak-sepele dan $x\in I\subseteq[x-\delta,x+\delta]$, berlaku
$|f(x)-\Bover1{\mu I}\int_If|\le\epsilon$.
%223D

\leader{223Y}{Latihan lanjutan (a)} Misalkan $E$, $F\subseteq\Bbb R$
himpunan-himpunan terukur, dan andaikan bahwa $0<\mu F<\infty$. Misalkan
$x\in\Bbb R$ sedemikian sehingga
$\lim_{h\downarrow 0}\Bover1{2h}\mu(E\cap[x-h,x+h])=1$. Tunjukkan bahwa

\Centerline{$\lim_{h\downarrow 0}\Bover{\mu(E\cap(x+hF))}{h\,\mu F}=1$.}

\noindent \Hint{terapkan 223Xg pada himpunan-himpunan berbentuk
$F\cap[-M,M]$.}
%223Xg 223B

\header{223Yb}{\bf (b)} Misalkan $\frak T$ keluarga himpunan-himpunan
terukur $G\subseteq\Bbb R$ sedemikian sehingga setiap titik dari $G$ merupakan
titik kerapatan dari $G$. {(i)} Tunjukkan bahwa $\frak T$ suatu topologi
pada $\Bbb R$. ({\it Petunjuk\/}: ambil $\Cal G\subseteq\frak T$. Menurut
215B(iv), terdapat suatu $\Cal G_0\subseteq\Cal G$ yang terhitung sedemikian
sehingga $\mu(G\setminus\bigcup\Cal G_0)=0$ untuk setiap $G\in\Cal G$.
Tunjukkan bahwa

\Centerline{$\bigcup\Cal G\subseteq\{x:\limsup_{h\downarrow 0}
  \Bover1{2h}\mu(\bigcup\Cal G_0\cap[x-h,x+h])>0\}$,}

\noindent sehingga
$\mu(\bigcup\Cal G\setminus\bigcup\Cal G_0)=0$.) {(ii)} Tunjukkan bahwa
suatu fungsi $f:\Bbb R\to\Bbb R$ terukur jika dan hanya jika fungsi itu
$\frak T$-kontinu pada hampir setiap $x\in\Bbb R$. ($\frak T$ adalah
{\bf topologi kerapatan} pada $\Bbb R$. Lihat 414P dalam Jilid 4.)
%223B

\spheader 223Yc Tunjukkan bahwa jika $f:[a,b]\to\Bbb R$ terbatas dan
kontinu terhadap topologi kerapatan pada $\Bbb R$, maka
$f(x)=\bover{d}{dx}\int_a^xf$ untuk setiap $x\in\ooint{a,b}$.
%223Yb 223B

\spheader 223Yd Tunjukkan bahwa suatu fungsi $f:\Bbb R\to\Bbb R$ kontinu
terhadap topologi kerapatan di $x\in\Bbb R$ jika dan hanya jika
$\lim_{h\downarrow 0}\bover1{2h}
  \mu^*\{y:y\in[x-h,x+h],\,|f(y)-f(x)|\ge\epsilon\}=0$ untuk
setiap $\epsilon>0$.
%223Yb 223B

\spheader 223Ye Suatu himpunan $A\subseteq\Bbb R$ disebut {\bf berpori}
di suatu titik $x\in\Bbb R$ jika
$\limsup_{y\to x}\Bover{\rho(y,A)}{|y-x|}>0$, dengan
$\rho(y,A)=\inf_{a\in A}|y-a|$. (Ambil $\rho(y,\emptyset)=\infty$.)
Tunjukkan bahwa jika $A$ berpori di setiap $x\in A$, maka $A$ terabaikan.
%223B

\spheader 223Yf Untuk suatu himpunan terukur $E\subseteq\Bbb R$, tuliskan
$\intstar E$ untuk himpunan titik-titik kerapatannya. Tunjukkan bahwa jika
$E$, $F\subseteq\Bbb R$ terukur, maka (i)
$\intstar(E\cap F)=\intstar E\cap\intstar F$
(ii) $\intstar E\subseteq\intstar F$ jika dan hanya jika
$\mu(E\setminus F)=0$ (iii) $\mu(E\symmdiff\intstar E)=0$ (iv)
$\intstar(\intstar E)=\intstar E$ (v) untuk setiap himpunan kompak
$K\subseteq\intstar E$, terdapat suatu himpunan kompak
$L\subseteq K\cup E$ sedemikian sehingga $K\subseteq\intstar L$.
%223B

\spheader 223Yg Misalkan $f$ suatu fungsi bernilai real terintegralkan yang
didefinisikan hampir di mana-mana dalam $\Bbb R$, dan $x$ sembarang titik
dari himpunan Lebesgue dari $f$. Tunjukkan bahwa untuk setiap $\epsilon>0$
terdapat $\delta>0$ sedemikian sehingga
$|f(x)\int g-\int f\times g|\le\epsilon\int g$ kapan pun
$g:\Bbb R\to\coint{0,\infty}$ sedemikian sehingga $g$ tak-menurun pada
$\ocint{-\infty,x}$, tak-menaik pada $\coint{x,\infty}$, dan nol di luar
$[x-\delta,x+\delta]$. \Hint{nyatakan $g$ sebagai limit hampir di mana-mana
dari fungsi-fungsi berbentuk
$\Bover{g(x)}{n+1}\sum_{i=0}^n\chi\ooint{a_i,b_i}$, dengan
$x-\delta\le a_0\le\ldots\le a_n\le x\le b_n\le\ldots\le b_0
\le x+\delta$.}
%223D

\spheader 223Yh Untuk setiap fungsi bernilai real terintegralkan $f$ yang
didefinisikan hampir di mana-mana dalam $\Bbb R$, misalkan $E_f$ adalah himpunan
Lebesgue dari $f$. Tunjukkan bahwa $E_f\cap E_g\subseteq E_{f+g}$ dan
$E_f\subseteq E_{|f|}$ untuk semua fungsi terintegralkan $f$, $g$.
%223D

\spheader 223Yi\dvAnew{2014}
Misalkan $E\subseteq\Bbb R$ suatu himpunan terukur tak terabaikan.
Tunjukkan bahwa $0$ termasuk dalam interior
$E-E=\{x-y:x$, $y\in E\}$.
%223B
}%end of exercises

\endnotes{
\Notesheader{223} Hasil-hasil dalam bagian ini dapat dipandang sebagai
pernyataan bahwa suatu fungsi terukur, dalam arti tertentu, `hampir kontinu';
bahkan 223Yb merupakan suatu upaya untuk memperjelas gagasan ini. Untuk fungsi
terintegralkan, kita mempunyai hasil-hasil yang lebih kuat, dan yang tampaknya
menjangkau paling jauh ialah 223D/223Ec.

Terdapat versi berdimensi $r$ dari semua teorema ini, dengan bola yang berpusat
di $x$ menggantikan interval $[x-h,x+h]$; saya memberikannya dalam 261C-261E.
%261C 261D 261E
Suatu gagasan baru diperlukan untuk versi berdimensi $r$ Teorema Kerapatan
Lebesgue (261C), tetapi selebihnya dapat digeneralisasi secara langsung. Suatu
perluasan yang kurang alami dan kurang penting, juga dalam \S261, melibatkan
fungsi-fungsi yang didefinisikan pada himpunan tak-terukur (bandingkan
223Xb-223Xe).
}%end of comment

\discrpage
